\section{Formulation}

\subsection{Constraint Systems in Hindley--Milner Inference}

Hindley--Milner type inference reduces to solving a finite system of
type-equality constraints generated by a syntax-directed traversal of the
program.
Formally, let $C = \{c_1, \ldots, c_m\}$ be a set of constraints of the form
$\tau_i = \tau_j$, where each $\tau_k$ is a type expression over a set of type
variables.
Each constraint $c_i$ carries a positive weight $w_i \in \mathbb{R}_{>0}$
reflecting the confidence assigned to that typing judgment~\cite{pavlinovic2014}.
The system $C$ is \emph{satisfiable} if there exists a substitution $\sigma$
such that $\sigma(\tau_i) = \sigma(\tau_j)$ for all $c_i \in C$; otherwise it
is \emph{unsatisfiable}.

\subsection{Minimal Correction Subsets and the MWMCS Problem}

When $C$ is unsatisfiable, a \emph{Minimal Correction Subset} is a
subset $M \subseteq C$ such that $C \setminus M$ is satisfiable, and no proper
subset of $M$ has this property~\cite{marques-silva2013}.
Each element of an MCS is a candidate error location.
In the example from Case Study~1, \texttt{let f x = x + 1 in f True}, the
constraints $c_0\colon \alpha = \texttt{Int}$ and
$c_1\colon \alpha = \texttt{Bool}$ form the sole MUS.
There are two minimum-weight correction subsets, $\{c_0\}$ and $\{c_1\}$,
each of total weight $1.0$.

The \emph{Minimum-Weight MCS} problem asks:
\begin{quote}
  Given an unsatisfiable weighted constraint system $(C, w)$, find a correction
  subset $M^* \subseteq C$ minimising $\sum_{c \in M^*} w_c$.
\end{quote}

By the hitting-set duality of~\cite{marques-silva2013}, every correction subset
is a hitting set of the family of Minimal Unsatisfiable Subsets of $C$, and
every hitting set of all MUSes is a correction subset.
MWMCS is therefore directly related to Minimum-Weight Hitting Set on the MUS
hypergraph.
The general hitting-set problem is NP-hard and is inapproximable within
$(1-\epsilon)\ln|C|$ for any $\epsilon > 0$ unless
$\mathrm{P} = \mathrm{NP}$~\cite{feige1998}.

\subsection{Research Question and Scope}

This paper investigates the following question empirically:

\begin{quote}
  \emph{On HM-shaped synthetic constraint instances, does a greedy
  hitting-set approximation achieve a ratio strictly greater than~$1$ relative
  to the exact MWMCS solution, and if not, what structural properties of the
  instances explain the absence of a gap?}
\end{quote}

We restrict scope to three synthetic generator families, described in
Section~\ref{sec:methodology}, spanning a single-conflict structure
Dataset~A, overlapping conflicts on a shared type variable Dataset~B, and
HM-shaped constructor mismatches Dataset~C.
For each family, we evaluate $|C| \in \{10, 20, 50, 100, 200\}$ with 20
instances per cell.
We use an exact MaxSMT-style baseline implemented with Z3 as the reference
optimum and do not claim a theoretical approximation bound for the lazy
implementation.
